\documentclass[11pt]{article}
\usepackage[margin=0.9in]{geometry}
\usepackage{amsmath,amssymb,amsthm,mathtools}
\usepackage[T1]{fontenc}
\usepackage[utf8]{inputenc}
\usepackage{enumitem}
\usepackage{booktabs,longtable,array,seqsplit}
\usepackage{hyperref}
\usepackage{xcolor}
\usepackage{microtype}
\hypersetup{colorlinks=true,linkcolor=blue!45!black,citecolor=blue!45!black,urlcolor=blue!45!black}
\setlength{\emergencystretch}{6em}

\newtheorem{theorem}{Theorem}[section]
\newtheorem{proposition}[theorem]{Proposition}
\newtheorem{lemma}[theorem]{Lemma}
\newtheorem{corollary}[theorem]{Corollary}
\newtheorem{claim}[theorem]{Claim}
\theoremstyle{definition}
\newtheorem{definition}[theorem]{Definition}
\newtheorem{assumption}[theorem]{Assumption}
\theoremstyle{remark}
\newtheorem{remark}[theorem]{Remark}

\newcommand{\R}{\mathbb{R}}
\newcommand{\C}{\mathbb{C}}
\newcommand{\N}{\mathbb{N}}
\newcommand{\Hsp}{\mathcal{H}}
\newcommand{\Fsp}{\mathcal{F}}
\newcommand{\Id}{\mathrm{Id}}
\newcommand{\tr}{\operatorname{tr}}
\newcommand{\diag}{\operatorname{diag}}
\newcommand{\rank}{\operatorname{rank}}
\newcommand{\spec}{\operatorname{spec}}
\newcommand{\spanop}{\operatorname{span}}
\newcommand{\dist}{\operatorname{dist}}
\newcommand{\range}{\operatorname{ran}}
\newcommand{\Ker}{\operatorname{ker}}
\newcommand{\Pinch}{\mathcal{C}}
\newcommand{\Off}{\widetilde{\mathcal{C}}}
\newcommand{\Pperp}{P^{\perp}}
\newcommand{\Qperp}{Q^{\perp}}
\newcommand{\op}{\mathrm{op}}
\newcommand{\F}{\mathrm{F}}
\newcommand{\HS}{\mathrm{HS}}
\newcommand{\ip}[2]{\left\langle #1,#2\right\rangle}
\newcommand{\norm}[1]{\left\lVert #1\right\rVert}
\newcommand{\abs}[1]{\left\lvert #1\right\rvert}
\newcommand{\code}[1]{\nolinkurl{#1}}

\newenvironment{argument}{\paragraph{Argument.}\begin{enumerate}[leftmargin=2em]}{\end{enumerate}}
\newenvironment{proofoutline}{\paragraph{Proof architecture.}\begin{enumerate}[leftmargin=2em]}{\end{enumerate}}
\newcommand{\sourceanchor}[1]{\par\smallskip\noindent\textbf{Source anchor.} #1\par\smallskip}
\newcommand{\sourcecomment}[1]{\begin{remark}#1\end{remark}}
\newenvironment{sourceguide}{\begin{description}[style=nextline,leftmargin=3.3cm,labelwidth=3.1cm]}{\end{description}}
\newenvironment{formalizationmap}{\begin{longtable}{@{}p{0.20\textwidth}p{0.31\textwidth}p{0.40\textwidth}@{}}\toprule Source item & Modern mathematical content & Repository status \\\midrule\endfirsthead\toprule Source item & Modern mathematical content & Repository status \\\midrule\endhead}{\bottomrule\end{longtable}}

\title{Halmos (1969), \emph{Two Subspaces}:\\canonical projection geometry for tangent maps and direct rotations}
\author{}
\date{Formalization-oriented distilled reconstruction and Davis comparison}
\begin{document}
\maketitle

\section*{Source map and scope}
\begin{sourceguide}
\item[Primary source] Paul R. Halmos, \emph{Two Subspaces}, Transactions of the American Mathematical Society 144 (1969), 381--389.
\item[Companion source] Chandler Davis, \emph{The Rotation of Eigenvectors by a Perturbation} (1963), especially the canonical rotation layer already reconstructed in \code{Davis1963_rotation_of_eigenvectors.tex}.
\item[Modern checks] Later formulations of Halmos's two-projection theorem and graph-subspace geometry; the repository's \code{PrincipalAngles.lean}, \code{DavisKahan/Experimental/InfiniteDimensional/Core/OperatorAngle.lean}, \code{DavisKahan/Experimental/FiniteDimensional/DirectRotation.lean}, and experimental \code{GraphSubspace.lean}.
\item[Purpose] Extract the exact decomposition and functional calculus needed for proof-carrying graph operators, principal tangent singular values, projection-distance criteria, and canonical direct rotation.
\end{sourceguide}

Halmos's theorem is the right geometric foundation because it separates all degeneracies before introducing an angle operator. In finite dimension this becomes a direct sum of trivial intersection blocks and generic principal two-planes. Almost every hard tangent or direct-rotation proof becomes scalar once this decomposition is available.

\section{The four intersection subspaces and generic remainder}

Let $P$ and $Q$ be orthogonal projections on a finite-dimensional Hilbert space $E$. Define
\[
  E_{11}:=\range P\cap\range Q,
  \qquad
  E_{10}:=\range P\cap\ker Q,
\]
\[
  E_{01}:=\ker P\cap\range Q,
  \qquad
  E_{00}:=\ker P\cap\ker Q.
\]
These subspaces are mutually orthogonal and reduce both projections. Let $E_g$ be their common orthogonal complement. Then
\[
  E=E_{11}\oplus E_{10}\oplus E_{01}\oplus E_{00}\oplus E_g.
\]
On the first four pieces, $(P,Q)$ is respectively
\[
  (I,I),\quad(I,0),\quad(0,I),\quad(0,0).
\]
The pair on $E_g$ is in generic position:
\[
  \range P\cap\range Q=0,
  \quad
  \range P\cap\ker Q=0,
\]
\[
  \ker P\cap\range Q=0,
  \quad
  \ker P\cap\ker Q=0.
\]

This decomposition must be explicit in a maximal finite-dimensional API. Totalized tangent definitions often hide $E_{10}$, the precise pole subspace where a principal angle equals $\pi/2$.

\section{Generic-position block model}

On the generic part, choose $K:=\range(P|_{E_g})$. There is a unitary identification
\[
  E_g\cong K\oplus K
\]
under which
\[
  P=\begin{bmatrix}I&0\\0&0\end{bmatrix}
\]
and
\[
  Q=\begin{bmatrix}
      C^2 & CS\\
      CS & S^2
    \end{bmatrix}.
\]
Here $C$ and $S$ are commuting positive contractions on $K$ satisfying
\[
  C^2+S^2=I,
\]
and generic position means
\[
  \ker C=0,
  \qquad
  \ker S=0.
\]
In finite dimension, $C$ and $S$ are simultaneously diagonalizable. Their eigenvalues are
\[
  \cos\theta_j,
  \qquad
  \sin\theta_j,
  \qquad 0<\theta_j<\frac\pi2.
\]
Thus the generic block is an orthogonal direct sum of ordinary two-dimensional principal planes.

\begin{remark}
The off-diagonal block is $CS$, not an independent operator. Proving commutation and the identity $C^2+S^2=I$ before diagonalization is what makes all later functional-calculus formulas basis independent.
\end{remark}

\section{Coordinate-free closeness, separation, and angle operators}

Halmos-style full-space operators are
\[
  \mathcal C(P,Q):=PQP+P^\perp Q^\perp P^\perp,
\]
\[
  \mathcal S(P,Q):=PQ^\perp P+P^\perp QP^\perp.
\]
They satisfy
\[
  0\le\mathcal C\le I,
  \qquad
  0\le\mathcal S\le I,
  \qquad
  \mathcal C+\mathcal S=I.
\]
Define the full operator angle by
\[
  \Theta(P,Q):=\arccos\sqrt{\mathcal C(P,Q)}.
\]
Then
\[
  \cos^2\Theta=\mathcal C,
  \qquad
  \sin^2\Theta=\mathcal S.
\]
This full-space definition is symmetric in the two subspaces and records unmatched dimensions. Restricted angle operators on $\range P$ are convenient for directed theorems, but they should be derived from the full geometry.

\section{Principal singular-value dictionary}

Let $U=\range P$ and $V=\range Q$. The singular values of the directed blocks are:
\[
  s(QP|_U)=\cos\theta_j,
\]
\[
  s(Q^\perp P|_U)=\sin\theta_j.
\]
For equal rank, the nonzero singular values of $P-Q$ are the principal sines, each with the standard doubled multiplicity on the generic planes. In particular,
\[
  \|P-Q\|=\max_j\sin\theta_j.
\]
The double-angle cross block is
\[
  2P^\perp QP,
\]
with singular values $\sin(2\theta_j)$.

The tangent graph map on the transverse part is
\[
  T=SC^{-1},
\]
with singular values
\[
  \tan\theta_j.
\]
This theorem is the exact missing dictionary behind \code{singularValues_graphOperator}. It should be proved once from simultaneous diagonalization, not by a separate matrix calculation for each tangent construction.

\section{Transversality and graph subspaces}

Let $U$ be fixed and let $V$ have the same finite dimension. The following are equivalent:
\begin{enumerate}
\item $V\cap U^\perp=\{0\}$;
\item the compression $P_U|_V:V\to U$ is invertible;
\item $V$ is the graph of a unique linear map $T:U\to U^\perp$;
\item every principal angle is $<\pi/2$;
\item $\|P_U-P_V\|<1$.
\end{enumerate}
The graph is
\[
  V=\{u+Tu:u\in U\}.
\]
Its normalized isometric embedding is
\[
  u\longmapsto (I+T^*T)^{-1/2}u
    +T(I+T^*T)^{-1/2}u.
\]
Hence
\[
  C=(I+T^*T)^{-1/2},
  \qquad
  S=T(I+T^*T)^{-1/2},
\]
and
\[
  SC^{-1}=T.
\]
Therefore the operator angle satisfies
\[
  \tan\Theta=|T|
\]
on the source coordinate space, and
\[
  \tan\theta_{\max}=\|T\|.
\]

\subsection{Unequal dimensions}

If $\dim V>\dim U$, transversality to $U^\perp$ is impossible. If $\dim V<\dim U$, $V$ may be the graph over a proper subspace of $U$. Thus the equal-rank graph theorem and the directed rectangular sine theorem have different maximal domains. An API should not use one implicit dimension convention for both.

\section{Reflection algebra and doubled angles}

Define the reflections
\[
  J_P:=2P-I,
  \qquad
  J_Q:=2Q-I.
\]
On a principal two-plane with angle $\theta$, the product $J_QJ_P$ is rotation by $2\theta$. Consequently,
\[
  \frac12\|J_Q-J_P\|
  =\|P-Q\|
  =\sin\theta_{\max},
\]
and the skew or off-diagonal block of the reflection product encodes $\sin(2\Theta)$.

This explains the Davis--Kahan reflection proof: conjugating an operator by $J_Q$ turns the angle between $P$ and $Q$ into a single-angle comparison between $P$ and $J_QP$, whose principal angles are doubled on the correct component.

The branch problem at $\pi/4$ is visible here. Rotation by $2\theta$ cannot distinguish $\theta$ from $\pi/2-\theta$ using sine alone. A tangent or inverse-cosine reconstruction must select the component containing zero.

\section{Canonical direct rotation}

On each generic principal plane define
\[
  W_\theta=
  \begin{bmatrix}
    \cos\theta&-\sin\theta\\
    \sin\theta&\cos\theta
  \end{bmatrix}.
\]
Extend it by the identity on $E_{11}\oplus E_{00}$. Defect subspaces $E_{10}$ and $E_{01}$ require equal dimension for a unitary intertwiner and destroy uniqueness unless additional data are supplied. For an acute equal-rank pair they are absent, and the direct rotation is canonical.

In generic block form,
\[
  W=
  \begin{bmatrix}
    C&-S\\
    S&C
  \end{bmatrix}.
\]
It satisfies
\[
  WPW^*=Q,
  \qquad
  W^2=J_QJ_P.
\]
If
\[
  K:=QP+Q^\perp P^\perp,
\]
then $K$ has block form
\[
  K=\begin{bmatrix}C^2&-CS\\CS&C^2\end{bmatrix},
\]
and the direct rotation is its polar factor on the acute part:
\[
  W=K|K|^{-1}.
\]
This is the cleanest construction for Lean because uniqueness of the polar factor packages basis independence.

\section{Minimality statements and their hypotheses}

The direct rotation is ``most direct'' in several inequivalent senses.

\begin{enumerate}
\item It minimizes Frobenius displacement among unitaries sending $U$ to $V$.
\item It minimizes the sum of squared angles moved by an orthonormal basis.
\item It minimizes maximal vector displacement.
\item Under an appropriate angle cap, positive displacement singular values are weakly majorized by those of any competing intertwiner, yielding UI-norm minimality.
\end{enumerate}

The last statement is the most delicate. On a principal plane, the direct rotation chooses the short arc. Without an acute or stronger angle cap, another unitary can change branches and invalidate a naive scalar monotonicity argument. The current \code{DirectRotation.lean} comments correctly warn that the UI-norm claims must retain the qualifications in Davis--Kahan Section 4.

\section{Lean-sized theorem seams}

A reusable finite implementation should prove the following in order.

\begin{enumerate}
\item Orthogonal decomposition into $E_{11},E_{10},E_{01},E_{00},E_g$.
\item Generic block equivalence $E_g\simeq K\oplus K$.
\item Positive commuting $C,S$ with $C^2+S^2=I$.
\item Simultaneous eigenbasis and principal-plane decomposition.
\item Singular values of cosine, sine, difference, and doubled cross blocks.
\item Equivalence of transversality, graph representation, and projection norm $<1$.
\item Graph embedding formulas and $\tan\Theta=|T|$.
\item Reflection product and doubled-angle formulas.
\item Polar-factor construction of direct rotation.
\item Uniqueness and intertwining.
\item Frobenius, operator-norm, and UI-norm minimality with exact angle hypotheses.
\end{enumerate}

\section{Repository map}

\begin{formalizationmap}
Intersection and generic decomposition & Four defect spaces plus principal generic part & Partially implicit in \code{PrincipalAngles.lean}; should become an explicit reusable theorem. \\
Cosine/sine blocks & Positive commuting pair and principal singular values & Much of the sine layer exists; the tangent singular-value dictionary remains incomplete. \\
Graph theorem & $\|P-Q\|<1$, transversality, unique angular operator & Experimental \code{GraphSubspace.lean} contains the intended statements with sorries. \\
Tangent identity & $\tan\Theta=|T|$ and maximal-angle norm formula & Central missing seam for \code{TanTheta.lean}. \\
Reflection geometry & Product of reflections and doubled angles & Used by \code{SinTwoTheta.lean}; should be centralized. \\
Direct rotation & Polar factor, trigonometric block, square identity, uniqueness & \code{DirectRotation.lean} has the architecture but several root proofs remain. \\
Minimality & Frobenius/UI/max displacement & Largest unfinished Davis geometry cluster. \\
\end{formalizationmap}

\section*{Bibliography}
\begin{thebibliography}{9}
\bibitem{Halmos1969} P. R. Halmos, \emph{Two Subspaces}, Trans. Amer. Math. Soc. 144 (1969), 381--389.
\bibitem{Davis1963} C. Davis, \emph{The Rotation of Eigenvectors by a Perturbation}, J. Math. Anal. Appl. 6 (1963), 159--173.
\bibitem{DK1970} C. Davis and W. M. Kahan, \emph{The Rotation of Eigenvectors by a Perturbation. III}, SIAM J. Numer. Anal. 7 (1970), 1--46.
\end{thebibliography}

\end{document}
